\chapter{Minimally Invasive Lumbar Spinal Fusion}

\section*{Overview}
Lumbar spinal fusion joins two or more vertebrae to stabilize the lower spine. Minimally invasive techniques use smaller incisions and specialized instruments.

\section*{Why It Is Done}
It may be considered for instability, degenerative spondylolisthesis, recurrent disc disease, deformity, or persistent pain and neurologic symptoms after nonoperative treatment.

\section*{How It Is Performed}
Under general anesthesia, the surgeon places screws and rods and inserts bone graft, often through tubular retractors or a lateral or posterior approach. Decompression may be performed at the same operation.

\section*{Preparation, Risks, and Recovery}
Imaging and medical optimization are required. Risks include infection, bleeding, nerve injury, dural tear, nonunion, hardware failure, adjacent-segment disease, and persistent pain. Recovery includes wound care, gradual activity, and rehabilitation.

\section*{References}
\begin{enumerate}
  \item MedlinePlus. Spinal Fusion. U.S. National Library of Medicine; 2025.
  \item Current Medical Diagnosis and Treatment. McGraw-Hill; 2025.
\end{enumerate}
\clearpage
